\documentclass[11pt,a4paper]{article}
\usepackage[utf8]{inputenc}
\usepackage[T1]{fontenc}
\usepackage{lmodern,amsmath,amssymb,amsthm,mathtools,booktabs,enumitem,microtype}
\usepackage{natbib}
\usepackage[margin=1in]{geometry}
\usepackage{xr-hyper}
\usepackage{hyperref}
\externaldocument[GE-]{01_law_of_graphic_equalisation_theory_referenced}
\externaldocument[IS-]{00_is}
\externaldocument[PH-]{01a_graphic_equalisation_physics_grounding}
\newtheorem{definition}{Definition}[section]
\theoremstyle{remark}\newtheorem{remark}[definition]{Remark}
\usepackage{fancyhdr}
\pagestyle{fancy}
\fancyhf{}
\fancyfoot[C]{\thepage}
\fancyfoot[R]{\scriptsize\textsc{draft-2026-09-02.1\,r10}}
\renewcommand{\headrulewidth}{0pt}
\fancypagestyle{plain}{\fancyhf{}\fancyfoot[C]{\thepage}%
  \fancyfoot[R]{\scriptsize\textsc{draft-2026-09-02.1\,r10}}\renewcommand{\headrulewidth}{0pt}}
\title{Frames, Relation, and the Empty Limit\\\large A Metaphysical Companion to Graphic Equalisation}
\author{Matthew Parslow\\Independent Researcher}
\date{\textbf{Draft \texorpdfstring{\textemdash}{---}\ draft-2026-09-02.1}\\[2pt]2 September 2026 \texorpdfstring{\textbullet}{\textperiodcentered}\ document revision 10}

\begin{document}
\maketitle
\begin{center}
\fbox{\parbox{0.92\textwidth}{\small
\textbf{DRAFT --- draft-2026-09-02.1 r10} (2 September 2026).
This document is under active revision and is circulated for comment, not as a
finished result. It is the metaphysical companion to \emph{The Law of Graphic
Equalisation --- Theory} (referred to below as Graphic Equalisation), and inherits
its definitions, epistemic boundary, and level-of-claims. The other documents
referred to below carry short names too, and those are the only forms used:
\emph{Is} (the base paper) and the physics companion (\emph{\ldots\ Physics
Grounding}). It is kept separate from Graphic Equalisation deliberately: the
material here is more speculative than that paper's, and should not borrow that
paper's standing.}}
\end{center}
\vspace{1em}

\begin{abstract}
Two questions are taken up that Graphic Equalisation raises and should not answer. First, whether a
well-formed frame approaches a view from nowhere, and in what sense the limit exists. Second, whether
relations are prior to relata, and what the substrate can say about a question of that shape. The
answers offered are deflationary in both cases and neither is claimed as new: the limit of
well-formedness is a purpose-relative optimum rather than a perspectiveless standpoint, and the
priority question is underdetermined by this substrate rather than settled or meaningless. What the
account contributes is a mechanism for a gradient that has usually been described without one.
\end{abstract}

\section{Scope}\label{sec:scope}
This paper depends on Graphic Equalisation for its vocabulary, and particularly on that paper's
definition of a frame as an instantiated boundary relation (Graphic Equalisation's
Definition~\ref{GE-def:frame}), its definition of well-formedness (its
Definition~\ref{GE-def:well-formed}), and its single prohibition with the formal analogue attached to
it (its Proposition~\ref{GE-prop:prohibition}). Nothing here is used by that paper; the dependency
runs one way. It is kept separate because a conclusion is not where speculative metaphysics should
acquire the standing of the work it concludes.

\section{Objectivity as a direction}\label{sec:direction}

What makes a frame a \emph{somewhere} rather than an \emph{anywhere} is its unnecessary structure:
what it holds that does no work, the purposes encoded in it that it does not serve, the distinctions
not earning their cost. Well-formedness is defined by the removal of exactly that, and that
definition's cost of held structure (Graphic Equalisation's Definition~\ref{GE-def:well-formed}) is
the mechanism this paper uses throughout; so a better-formed frame is a less particular one, and
frames doing the same work grow more alike as each is better formed (Graphic Equalisation's
Proposition~\ref{GE-prop:convergence-observed}, which rests on the near-uniqueness posit of its
Hypothesis~\ref{GE-hyp:cheap-cut}). Objectivity on this reading is a direction rather than a standpoint, and the direction is
real and traversable, which is why enquiry makes progress that feels like approaching something.

It is tempting to call the limit of that direction a view from nowhere. The temptation should be resisted, and the reason is internal rather
than modest. Well-formedness is relative to \emph{what work} a frame is doing, and work is a purpose,
and a purpose is a somewhere. Convergence under better-formedness is therefore convergence toward a
\emph{shared} somewhere --- the optimum for that work --- and not toward a perspectiveless standpoint.
Two frames doing the same work converge on each other. Two frames doing different work do not, and
nothing in the mechanism suggests they should.

Three further obstacles stand between the direction and any terminus, and they are independent. The
first two are Graphic Equalisation's, being two of the three clauses of its single prohibition
(Proposition~\ref{GE-prop:prohibition}): no level is self-transparent, so even a perfectly
well-formed frame cannot observe its own running; and a frame's measurement of its own throughput
returns its own unit whatever the underlying rate. The third is this paper's own addition. The
last of the particularity cannot be stripped without removing the criterion of well-formedness along
with it, since that criterion is the work: strip the purpose and there is no ranking left, hence no
cut, hence no frame.

\begin{remark}[The limit point is empty]\label{rem:empty-limit}
The set of frames with no particularity has no members, and the only object answering to the
description --- nothing held, no work done, no purpose, hence no cut --- is the empty set, which is
not a frame and does nothing. Approaching the limit is genuine progress; arriving would be
dissolution rather than arrival.
\end{remark}

This is the empty end of \emph{Is}'s two zeros (\emph{Is}, Section~\ref{IS-sec:one-not-enough}),
reached from the direction of well-formedness rather than of relation; Graphic Equalisation's
zero-bits reading (Graphic Equalisation's Remark~\ref{GE-rem:everything-nothing}) is the same point measured.

Two things that remark does not license. It is not a claim that objectivity is unavailable or that
enquiry is futile, since the direction is real and the convergence is measurable. And it is not an
appeal to the uncomputability of a shortest description, which would not support it: minimal
descriptions are not unique, description length is machine-relative up to a constant --- an objection
Graphic Equalisation accepts as a thesis rather than deflects, the constant being the frame (Graphic
Equalisation's Remark~\ref{GE-rem:mdl-ontic}) --- and the convergence claimed here is relative to what supplies a frame's primitives, for the same reason
everything else in this account is.

\begin{remark}[Prior positions on the empty limit]\label{rem:nowhere-prior}
Neither the graded reading of objectivity nor the suspicion of its limit is new here, and the nearest
sources should be described accurately rather than enlisted. \citet{nagel1986} holds objectivity to be
achieved by degrees of stepping back, which is the graded half; but his limit is \emph{contentful} and
his residual is a completeness residual --- the objective conception omits the subjective, which is
real --- so he is not a predecessor of the claim that the limit is empty, and is more often read as
its target. The absolute conception, a limiting notion maximally free of perspective, is proposed by
\citet{williams1978}, whose own reservations concern grounding ethics in it rather than its
coherence; the suspicion that no such conception is available belongs to his critics rather than to
him. What is offered here is neither position but a mechanism for the gradient both describe: the
cost of what a frame holds that does no work.
\end{remark}

\section{Relations and relata}\label{sec:relation}

\emph{Is} holds that to be is to stand in relation (\emph{Is}, Proposition~\ref{IS-prop:is}) and that
to exist is to cohere in it (\emph{Is}, Proposition~\ref{IS-prop:cohere}), and Graphic Equalisation
that rules running is what a world in which anything \emph{happens} costs, not the condition of there
being anything at all (Graphic Equalisation's Remark~\ref{GE-rem:only-g}). What is added here is
confined to the \emph{priority} question, which neither settles. Particularity is relation --- to be cut out is to be
related to what one is cut from --- so the imperfection of Section~\ref{sec:direction} is not a defect
accompanying existence but the relation constituting it.

The priority question follows immediately and should not be answered. Are relations prior to relata,
or relata to relations?

\begin{remark}[Underdetermined by this substrate]\label{rem:underdetermined}
To determine priority one needs a description of one side independent of the other. Within this
account there is none, because a relatum's individuation is the drawing of a boundary and the
boundary is the relation. But that is a fact about \emph{this} substrate's resources and not a
demonstration that the question has no answer: the relata-first theorist denies precisely the premise
the argument runs on, so the argument presupposes one horn and cannot settle between them. The honest
statement is that the question is \emph{underdetermined by this substrate}, not that it is unknowable
\emph{simpliciter}. The stronger form does not survive the observation that its premise is the position it purports to leave open.
\end{remark}

What can be said beyond the underdetermination is what Graphic Equalisation says of any such choice.
Relation-first is the better-formed hypothesis here: it requires no posit this account does not
already use, where relata-first requires a grounding claim that does no work in it. Better-formedness
licenses assuming it and does not make it true, which is the same form of answer the physics
companion gives to the choice of a synchronisation convention (the physics companion's
Remark~\ref{PH-rem:oneway-rule}, the better-formed-hence-assumed argument, with its
Remark~\ref{PH-rem:reichenbach-malament} for the two halves of that situation) --- and the recurrence
of that form across a physical and a metaphysical question is the most interesting thing in this
paper.

That horn is one \emph{Is} has already taken on its own argument: it holds relational determinacy to
be a commitment rather than a result, and the one place to disagree with it (\emph{Is},
Section~\ref{IS-sec:claim}), and holds the \emph{constructive} priority of relata unavailable at that
level, an order of construction being a sequence and there being no time there for one to run in
(\emph{Is}, Section~\ref{IS-sec:prior}); what is claimed here is the underdetermination of the
\emph{metaphysical} priority, which is compatible with both and does not inherit either.

\begin{remark}[Prior positions on relation and relata]\label{rem:relation-prior}
The positions and their characterisations --- \citet{ladymanross2007} on the strong horn,
\citet{esfeld2004} on the moderate one, and the far older denial that anything possesses intrinsic,
self-establishing existence \citep{nagarjuna1995} --- are set out in \emph{Is}
(Section~\ref{IS-sec:prior}) and are not restated here; what this section adds is the qualification
that Madhyamaka declines the positive thesis as well, so it is an ally for the negative half of this
section and not for the positive one.
\end{remark}

\begin{remark}[What the Yoneda lemma does and does not supply]\label{rem:yoneda}
It is tempting to enlist a mathematical statement here, and the temptation should be qualified rather
than indulged. The Yoneda lemma, and more precisely the corollary that the Yoneda embedding is full
and faithful, gives that an object of a locally small category is determined up to isomorphism by the
\emph{functor} of morphisms into it, with the isomorphism natural \citep[\S III.2]{maclane1971}. Three
things follow for the use made of it. Naturality is not decoration: bare bijections between hom-sets
do not determine the object, so the loose gloss ``an object is its relations'' is strictly weaker
than the lemma and false as stated. Morphisms are directed and composable and a category is not a
relational structure, so ``relation'' is not a faithful translation. And the result is internal to a
fixed ambient category, with the object given in advance --- it is not a metaphysical thesis about
what there is. What it does supply is a precise sense in which relational data suffices to
individuate within a settled context, which is suggestive of the position of this section and is not
an argument for it.
\end{remark}

\section{A worked instance}\label{sec:chicken}

The familiar case is worth taking up, both because it locates the material and because the standard resolution is stronger than it first appears. Asked which came first, the chicken or the egg, the temporal question has an answer and vagueness
does not block it. Under every admissible sharpening of ``chicken'', the genetic novelty is in the
zygote, so the first organism satisfying the sharpening developed inside an egg laid by a bird that
did not. The answer is invariant under precisification, which is exactly the case in which vagueness
is harmless, and it does not require a sharp boundary to exist.

Where the puzzle's force actually lives is an ambiguity in \emph{egg} rather than in \emph{chicken}:
an egg laid by a chicken and an egg containing a chicken are different objects, and the two readings
give opposite answers. That is a disambiguation and not a discovery about carving.

What the substrate contributes is confined to the second question the puzzle can be asked as: whether
the kind is prior to its instances, or the instances to the kind. That is Section~\ref{sec:relation}'s
question and is underdetermined for the reason given there. The diagnosis worth keeping is of the
puzzle's \emph{feel} rather than of its content --- pressed on the temporal reading one retreats to
the constitutive, and pressed on the constitutive one retreats to the temporal, so an answer seems
always to be available somewhere other than where one is standing.

\section{Falsifiability}\label{sec:falsifiability}

\noindent\emph{On the tags.} Each commitment below is marked \textsc{[shape]} or \textsc{[detail]}.
A \textsc{[detail]} failing is a revision: the account survives with that mechanism replaced. A
\textsc{[shape]} failing costs \emph{range} rather than a mechanism, because what fails is something
the substrate is claimed to require or to permit --- so the account does not hold where it claimed to,
which is the more serious of the two and is still a boundary rather than an annihilation. The marks
were assigned in a single pass under the rule just stated, and are the least-audited judgement in this
document --- a reader who thinks one is marked in the flattering direction is disagreeing with a fast
call, not an argued position.

This paper's claims are weaker than Graphic Equalisation's and its commitments are correspondingly
narrower.

\begin{enumerate}[leftmargin=*]
\item \textsc{[shape]}~\textbf{Convergence is purpose-relative.} Fails if frames doing demonstrably different work
converge on the same structure as each becomes better formed, which would indicate a
purpose-independent optimum and restore the terminus this paper denies. The convergence relied on is
Graphic Equalisation's (Proposition~\ref{GE-prop:convergence-observed}), so this item also inherits
that paper's near-uniqueness posit (Hypothesis~\ref{GE-hyp:cheap-cut}) and fails with it.
\item \textsc{[shape]}~\textbf{The limit is empty.} Fails if a frame can be exhibited that holds no structure not
earning its cost and still does work --- that is, if the criterion of well-formedness survives the
removal of the purpose it is relative to.
\item \textsc{[shape]}~\textbf{Underdetermination, not unknowability.} Fails if a discriminating observation between
relation-first and relata-first can be specified from within any frame, in which case the question is
settled rather than underdetermined.
\item \textsc{[detail]}~\textbf{The recurrence is one form of answer, not two.} Fails if the licence
invoked here for relation-first and the one invoked for a synchronisation convention in the physics
companion (that paper's Remark~\ref{PH-rem:oneway-rule}) turn out to rest on different grounds ---
one on cost and the other on something else --- in which case the recurrence is a resemblance and the
claim made for it above is empty.
\end{enumerate}

What would strengthen rather than falsify item 3: a demonstration that the argument of
Remark~\ref{rem:underdetermined} does not presuppose its conclusion, which would recover the stronger
unknowability claim rather than refute the weaker one made here --- a supersession and not a
falsification, which is why it is not in the list. And Remark~\ref{rem:yoneda}'s negative --- that
``an object is its relations'' is false as stated --- carries no exposure and needs none, being a
reading of a theorem rather than a commitment of this account: what could go wrong with it is a
misreading, which is what the remark is for.

\begin{thebibliography}{99}

\bibitem[Esfeld(2004)]{esfeld2004}
Michael Esfeld.
\newblock Quantum entanglement and a metaphysics of relations.
\newblock \emph{Studies in History and Philosophy of Modern Physics}, 35(4):601--617, 2004.

\bibitem[Ladyman and Ross(2007)]{ladymanross2007}
James Ladyman and Don Ross.
\newblock \emph{Every Thing Must Go: Metaphysics Naturalized}.
\newblock Oxford University Press, Oxford, 2007.

\bibitem[Mac Lane(1971)]{maclane1971}
Saunders Mac Lane.
\newblock \emph{Categories for the Working Mathematician}.
\newblock Springer, New York, 1971.

\bibitem[N\=ag\=arjuna(1995)]{nagarjuna1995}
N\=ag\=arjuna.
\newblock \emph{The Fundamental Wisdom of the Middle Way: N\=ag\=arjuna's M\=ulamadhyamakak\=arik\=a}.
\newblock Translated by Jay L. Garfield. Oxford University Press, New York, 1995.

\bibitem[Nagel(1986)]{nagel1986}
Thomas Nagel.
\newblock \emph{The View from Nowhere}.
\newblock Oxford University Press, New York, 1986.

\bibitem[Williams(1978)]{williams1978}
Bernard Williams.
\newblock \emph{Descartes: The Project of Pure Enquiry}.
\newblock Penguin, Harmondsworth, 1978.

\end{thebibliography}
\end{document}
